The final step at a single edge is to place the detecting
functional in the target lattice of
Proposition~\ref{prop:residue-plane-neighbor-test}.
The dual-plane substitution accounts for the vanishing at
the other incoming edges.
\paragraph{Dual planes at the target.}
\label{par:plane-dual-edge}
\hypertarget{source-plane-edge-covector}{}
For clarity write the distinguished edge in adapted coordinates as
\(L_-=b W^k\), \(L_+=b D_p^{-1}W^k\), where
\(D_t=\operatorname{diag}(tI_r,I_{k-r})\). Its dual edge goes
from \(L_+^\vee=b^{-t}D_p W^k\) to
\(L_-^\vee=b^{-t}W^k\).

At each Frobenius position the middle
frame on this dual edge is the inverse transpose of the original
middle frame: in the original frame use \(bD_w^{-1}\), and in
the dual use
\[
 b^{-t}D_p D_w^{-1}=b^{-t}D_w,
\]
with the
indicated Frobenius action on \(b\). Distinct incoming edges at
\(L_+\) give distinct outgoing edges at \(L_+^\vee\) by double
duality. There they have distinct residue planes, since a lattice
between a lattice and its \(p^{-1}\)-multiple is determined by
its image in that quotient. For the standard target,
\(bD_p^{-1}\in\mathrm{GL}_k(W)\); it gives the Witt frame
whose inverse transpose is \(b^{-t}D_p\). Thus a single choice
of distinguished frame in the one-sided realization handles all
incoming competitors.

\begin{lemma}[A one-sided plane symbol gives an integral dual-edge test]\label{lem:plane-dual-edge-test}
For the edge from \(bW^k\) to \(bD_p^{-1}W^k\), use
\(bD_w^{-1}\) and \(b^{-t}D_w\) as the respective middle frames.
Let a detecting
symbol of order \(o\) be realized to precision \(o+c_P\)
at the distinguished dual edge and at least that order
at its competitors. If \(D_X<c_P/2\), the resulting flat covector on the original edge has order \(o'\), conductor \(c'\) and depth \(D'\) satisfying
\[
 o'\le o,\qquad c'\ge c_P,\qquad D'\le D_X.
\]

\end{lemma}

\begin{proof}
\proofstep{The dual pairing on the target plane.}
In flat polynomial coordinates use the bilinear form
\(p^{\kappa_{\rm fl}}[P,Q]/\prod_i d_i!\).
It is integral on coefficient lattices in inverse transpose
frames. Indeed the factorial deficiency gives this in the
standard coordinates; the substitution identity
\[
 [P(MX),Q(X)]=[P(X),Q(M^tX)]
\]
follows by writing the
pairing as a constant-coefficient differentiation at zero.
The identity applies to each letter, also when it carries the
inverse transpose representation. The scalar adjustments to the two edge substitutions give 
\[
 w^{A_\mu}w^{A_\mu}=w^{2A_\mu}=p^{A_\mu},
\]
 precisely
the normalized edge pairing for the flat weight \(\mu\), by
Section~\ref{par:plane-middle-weights}.

Take the integral polynomial realizing a detecting symbol
\(\sigma\) of order \(o\), and view it as a flat section in the
dual vertex frame. Divide it by \(w^o\), and pair by this form to
obtain \(T\). At the target its order is at most \(o\), by
integrality before division. In normalized coordinates on the
dual edge its polynomial after division is \(\sigma+w^{c_P}E\)
with integral \(E\); on any other dual edge at the same source
it belongs to \(w^{c_P}O[\mathbf X]\). The inverse transpose
calculation proves \(T\in R_e^{\prime\mu}\) and
\[
 T\in w^{c_P}R_{e'}^{\prime\mu}
\]
on every competing incoming
edge. For any normalized integral input, the model error pairs into
\(w^{c_P}O\), because the factorial rescaling is integral on
the paired edge lattices. Its valuation is at least
\[
 c_P/2,
\]

strictly larger than the detected valuation bounded by \(D_X\).
The valuation of the actual pairing therefore equals that
of the model detecting value.

\end{proof}

\paragraph{Operators for the bounded remaining weights.}
\label{par:plane-rho}
First retain only the top positions in the full weight \(\lambda\).
Write \(\lambda=\mu+\eta\) there, with \(\mu\) the flat weight
and \(\eta\) the remaining weight. Its spread at each top is at
most \(c_0\). We need two facts about the Schur lattices
of Section~\ref{sec:edge-contraction}.

First, for each of the weights \(\xi=\lambda,\mu\) there is a
contravariant identification \(\Psi_\xi:V_\xi^*\longrightarrow
V_\xi\) satisfying, simultaneously in every vertex frame \(g\),
\[
 p^{h_\xi} H_\xi(g^{-t})\ \subseteq \Psi_\xi\bigl(H_\xi(g)^*\bigr)\ \subseteq\ H_\xi(g^{-t}).
\]
In this display \(g\) denotes the tuple of the component frames
and \(H^*\) the integral dual. The same inclusions hold with
\(R_e^{\prime\xi},R_{e^\vee}^\xi\) in the middle and outside,
since the two normalizations multiply the unscaled lattices by
the same \(w^{-A_\xi}\).
Here \(h_\xi\) is the sum over positions of
\(\binom k2\lfloor\log_p(l+k-1)\rfloor\) at positive spreads
\(l\), with zero at spread zero.

\begin{proposition}[Uniform operators for the remaining weights]\label{prop:plane-bounded-weight-operators}
With the bounded top and off-top spreads, the maps
\(\rho_\beta\) of Proposition~\ref{prop:plane-restore-covector}
satisfy its three lattice conditions with
\(s=rNc_0+2rn'\mathrm{cpp}\). Every primitive vector
has an image outside \(w^{2\kappa+1}R_e^\mu\) for some
such map. On the restored model family in the prescribed range of parameters, the target depth \(D\), order \(o\) and conductor \(c\) satisfy
\[
 D\le D_X+\kappa\le54Nk^2\log_p a,\qquad o\le t_F+s,\qquad c=c_\star.
\]

\end{proposition}

\begin{proof}
\proofstep{Normalizing one position in the tensor.}
To verify the normalization at one position, remove the determinant
twist. Take the matrix-polynomial Fock form divided by the
squared Fock norm of the highest polynomial \(v\), the product
of the initial principal minors prescribed by the shape.
This norm is a positive integer. Pairing with \(v\) gives evaluation
at the identity: both vanish outside its weight by row and column
degrees, and both take value one on \(v\). The highest weight
space has dimension one. For a divided-power word applied to
\(v\), transpose the word in the pairing. It follows that its
pairing with any integral polynomial in the Schur lattice is
integral.

Conversely, integrality against the Schur lattice forces
\[
 Q\in O[\mathbf X]
\]
in the matrix coordinates.
Indeed substitute a generic matrix \(M\) in the highest polynomial
by the row action. The coefficients of the resulting polynomial
in the entries of \(M\) are integral elements of the Schur
lattice, by the minor expansion. Pair it with \(Q\) and transpose
the row action. The result is \(Q\) evaluated on the generic
transposed matrix, by the identity evaluation just proved.
Hence every coefficient of \(Q\) is integral, and saturation
of the standard monomial lattice gives the assertion.
The same argument proves nondegeneracy. The divided-power span of
\(v\) contains
\[
 p^{\binom k2\lfloor\log_p(l+k-1)\rfloor}
\]

times the Schur lattice by the trace argument of
Section~\ref{sec:edge-contraction}. Thus the integral dual for the
normalized form lies between the two asserted lattices. Restoring
the determinant twists on inverse transpose frames cancels their
two scalar factors. Tensoring the forms and integral dual bases
adds the exponents and gives \(\Psi_\xi\). Its inverse pairing
is a nonzero scalar times the matrix-polynomial Fock form.
Here \(l+k-1\le a^2\) for both weights at a sufficiently late
host: the flat degree is at most \(2a\), the residual spread
is at most \(c_0\), and the host size bound gives
\(2a+c_0+k\le a^2\). Consequently
\[
 \kappa=h_\lambda+h_\mu\le4L.
\]

Second, for residual sections \(s,\beta\) consider the flat
endomorphism \(U_{s,\beta}\) defined by
\[
  [\theta(U_{s,\beta}F),\theta H]_{\rm mat}
       =[\theta F\cdot s,\theta H\cdot\beta]_{\rm mat}.
\]
The matrix-polynomial form restricted to the image of the flat
model is nondegenerate, by the preceding argument. The flat
insertion \(\theta\) uses the first column or the cofactors.
The maps \(U_{s,\beta}\) are linear in either residual input.
At a single position their span is stable under commutators
with the row operators \(E_{ij}=X_i\partial_j\).
In fact transpose adjunction and the product rule move the
commutator to the two terms applying row operators to \(s\) and
to \(\beta\), with opposite signs and transposed indices.

For cofactors the induced row action on the flat coordinates is
the negative transpose action plus the scalar trace term, by
Jacobi's identity. If the induced operator is \(A+\lambda I\), its commutator on
an endomorphism \(U\) is
\[
 (A+\lambda I)U-U(A+\lambda I)=AU-UA .
\]
Thus the determinant twist introduces no extra row weight.
The residual space is generated by products of flag minors with
at most \(c_0\) columns. Each row occurs at most once in a
column. Thus the commutator weight of a generator
\(U_{s,\beta}\) has each coordinate between \(-c_0\) and \(c_0\):
it is the difference, in one of the two orders, of the two
residual row weights. The determinant twists cancel in this difference.

\proofstep{Operators from the diagonal action.}
For degree \(d\) the endomorphism representation of
\(\operatorname{Sym}^d\) in characteristic zero has one summand for each of the highest weights
\[
 \lambda^{(i)}=(i,0,\ldots,0,-i),\qquad 0\le i\le d.
\]
Indeed \(\operatorname{Sym}^d\otimes
(\operatorname{Sym}^d)^*\) is the bidegree \((d,d)\)
space decomposed by the powers of the contraction in
Section~\ref{sec:edge-transfer}; its harmonic quotients are
the irreducible incidence spaces described there. For summand
\(i\) the row word \(E_{1k}^i\) has that highest weight and
is killed under commutator by every raising operator. Its action on a single monomial gives
\[
 E_{1k}^i X_k^d
 =\frac{d!}{(d-i)!}X_1^i X_k^{\,d-i}\ne0
 \qquad(0\le i\le d),
\]
including the identity at \(i=0\). Words of length at
most \(i\) form a commutator-stable span by the product rule and
the matrix-unit commutator. They therefore contain the entire
summand.

Complete reducibility now implies that the span of the
\(U_{s,\beta}\) lies in words of length at most \(c_0\). Indeed
any summand present has its highest weight present, and the
weight bound just proved gives \(i\le c_0\).
At several positions expand the residual sections in pure
tensors and use factorization of the Fock form. Each term gives
at most \(c_0\) row operators at each top. Changes of flat frame
preserve this word space by conjugation. The word bound also holds when the flat degree is smaller than the allowed word length.

For \(\beta\in R_{e^\vee}^{\eta}\), write \(m_\beta\) for
multiplication by that residual section. Define
\[
   \Lambda_\beta T
       =p^{h_\lambda}\Psi_\lambda^{-1}
                      m_\beta\Psi_\mu T,\qquad
   T(\rho_\beta u)=p^{A_\eta}(\Lambda_\beta T)(u).
\]
The second equality defines a linear map by finite-dimensional
duality. The edge normalizations give
\[
 b_e^\mu(\rho_\beta u,T)=b_e^\lambda(u,\Lambda_\beta T).
\]
For \(T\in R_e^{\prime\mu}\), multiplication on the dual edge and
the inclusions for \(\Psi_\xi\) give
\[
 \Psi_\mu T\in R_{e^\vee}^\mu,\qquad
 m_\beta\Psi_\mu T\in R_{e^\vee}^\lambda,\qquad
 p^{h_\lambda}R_{e^\vee}^\lambda\subseteq
                              \Psi_\lambda(R_e^{\prime\lambda}).
\]
Thus \(\Lambda_\beta T\in R_e^{\prime\lambda}\).
In unscaled dual bases the two lattice factors \(w^{-A}\) cancel
the pairing factor \(w^{2A}\). The preceding adjoint identity
and this perfectness imply
\[
 \rho_\beta R_e^\lambda\subseteq R_e^\mu.
\]

On an input
\(u=\theta F\cdot s\), the definition of the two forms shows
that its flat output is a scalar multiple of
\(U_{s,\beta}F\) in polynomial coordinates. This proves
membership in the exact operator span.

\proofstep{Reservations for the remaining weights.}
Put \(s_T=rNc_0\). Multiplying any \(\beta\) from the chosen
dual edge by \(w^{s_T}\) puts it in the dual-edge source vertex
lattice. To see this, at each position the edge normalization
subtracts the top \(r\)-row weight sum. The largest possible
difference between that sum and a sum on any
\[
 r
\]
rows is at
most \(r\) times the weight spread, by the tableau weight bound.
This proves the claim in an adapted diagonal frame, including
determinant twists, and hence in every frame. Every outgoing normalized dual edge contains the vertex lattice.
Hence the cost outside the distinguished edge is at most \(s_T\).
For the target vertex, test against \(T\in(H_+^\mu)^*\). The vertex
inclusions and \(p=w^2\) give
\begin{equation}\label{eq:edge-target-adjoint}
 \begin{aligned}
 w^{s_T}\Lambda_\beta T &\in(H_+^\lambda)^*,\\
 T(\rho_\beta u)&=w^{2A_\eta-s_T}
          (w^{s_T}\Lambda_\beta T)(u)\in w^{2A_\eta-s_T}O
                      \qquad(u\in H_+^\lambda).
 \end{aligned}
\end{equation}
Testing every integral dual functional in \eqref{eq:edge-target-adjoint} yields
\(\rho_\beta H_+^\lambda\subseteq w^{2A_\eta-s_T}H_+^\mu\).

There is also a lower bound for this family of maps. The full
lattice on the dual edge is the integral span of products
\(y\beta\) from the flat and residual edge lattices. In a standard
tableau product, the column heights of the sum shape are the
union, with multiplicities, of those of the two shapes. Split
the columns accordingly and take a tensor of these factorizations
over the positions. The frame action, determinant twist and
edge powers all multiply. Thus for every full dual integral
functional \(G\), \(\Psi_\lambda G\) is an integral sum of these
products.

If
\[
 \Psi_\lambda G=\sum_i y_i\beta_i,
\]
take
\(T_i=\Psi_\mu^{-1}(p^{h_\mu}y_i)\), which is flat
dual integral. The multiplier in \(\Lambda_\beta\) then gives
\[
 p^\kappa G=\sum_i\Lambda_{\beta_i} T_i
\]
with
\(\kappa=h_\lambda+h_\mu\). For the normalized pairings,
\[
 p^\kappa b_e^\lambda(u,G)
       =\sum_i b_e^\mu(\rho_{\beta_i}u,T_i).
\]
Suppose \(\rho_\beta u\in w^{2\kappa+1}R_e^\mu\) for every
integral choice of \(\beta\). Testing the displayed sum against the
integral flat duals then gives, for every \(G\in R_e^{\prime\lambda}\),
\[
 p^\kappa b_e^\lambda(u,G)\in w^{2\kappa+1}O
 \quad\Longrightarrow\quad b_e^\lambda(u,G)\in wO
 \quad(p^\kappa=w^{2\kappa}).
\]
Perfectness gives \(u\in w R_e^\lambda\), a contradiction for a
primitive \(u\). Thus at least one residual test detects it.

\proofstep{Restoring the non-top tensor factor.}
Restore the non-top tensor factor, of weight \(\xi\), by taking
all contractions on it with \(p^{A_\xi}G\),
\(G\in R_e^{\prime\xi}\), before applying the top maps.
In a pair of normalized dual bases these contractions are
ordinary coordinate extractions. Some extraction of a primitive
tensor remains primitive. Thus this step adds no depth. The
spread off the tops is at most \(2\mathrm{cpp}\), so
\[
 w^{2rn'\mathrm{cpp}}G
\]
belongs to the target dual vertex
lattice and hence to all competing incoming dual edge lattices.

For this inclusion, the passage from an intermediate to a target
frame expands the first \(r\) directions once more by \(w^{-1}\);
a coordinate functional then loses at most the difference of two
\(r\)-row weight sums. On the target itself the contraction
carries the additional scalar \(p^{A_\xi}\). Adding the competing-edge loss from these positions to the
top contribution gives the budget
\[
 s=s_T+2rn'\mathrm{cpp},\qquad
 2A_\eta+2A_\xi-s=2(A_\lambda-A_\mu)-s .
\]
Here the right-hand exponent is the resulting target-vertex exponent.
The flat
degree is zero at every omitted position, so this contraction
retains the top word allowance.

Apply the adjoint calculation to the primitive restored model
combination. The three resulting bounds are
\begin{equation}\label{eq:restored-covector-budget}
 \begin{aligned}
 \text{endpoint order}&\le t_F+s\le\lceil t_{\rm II}\rceil,\\
 \text{depth}&\le D_X+\kappa\le54Nk^2\log_p a\le D_{\rm II},\\
 \text{competing conductor}&\ge c_P-s=c_\star,
                    \qquad c_\star=8(D_{\rm II}+1).
 \end{aligned}
\end{equation}
The last depth comparison and the endpoint bound in
\eqref{eq:restored-covector-budget} are the host choices of Section~\ref{sec:host-edge-tests},
with its absolute coefficient taken sufficiently large.
On the source side cleaning
the normal-degree components of the original flat generators
costs at most \(\lfloor X\rfloor\) at the source, by the
one-sided lift and division by the order of each component.
Multiplication by a remaining-weight section from the chosen
edge costs at most \(s\) at that vertex, by the same spread
estimate. Thus the source order also fits. With \(c'=c_P\)
one has \[
 c'\ge2(D_X+\kappa+1),\qquad c'\ge c_\star+s.
\]
Section~\ref{par:plane-middle-weights} transfers the pairing
to the clean combination and gives its competing source
conductors. This completes the single-edge estimate with the
order and depth fixed on the host and the stated quantifier over
primitive combinations. \label{cor:plane-edge-covector-1}
\end{proof}
